\section{Introduction}
\label{sec:intro}

Contemporary scientific research is experiencing an unprecedented surge in publication velocity. In fields such as artificial intelligence, machine learning, and computational linguistics, hundreds of manuscripts are deposited to pre-print repositories like arXiv every week. For active investigators, maintaining an up-to-date mental model of current methodologies, benchmark results, and related literature has transformed from a routine scholarly task into a significant cognitive burden.

Historically, researchers have relied on keyword-based alerting systems (e.g., Google Scholar Alerts, arXiv RSS feeds) to monitor literature. However, lexical keyword matching suffers from fundamental information-retrieval deficiencies~\cite{manning2008introduction,robertson1976relevance}, notably synonymy (novel concepts expressed with unfamiliar terminology) and polysemy (identical terms possessing distinct technical definitions across subfields). Consequently, keyword feeds inevitably degrade into high-volume, noisy message streams that researchers routinely abandon due to alert fatigue.

Recent advancements in dense semantic representations~\cite{reimers2019sentencebert,beltagy2019scibert,cohan2020specter} and Large Language Models (LLMs)~\cite{vaswani2017attention} present a compelling alternative: encoding publication abstracts into semantic vector spaces to compute semantic similarity against a user profile, and delegating reading and synthesis to LLMs. Nevertheless, naive ``agentic'' literature monitoring architectures---which query generative LLMs to evaluate, score, and summarize every incoming paper---suffer from acute economic and computational limitations. Processing hundreds of daily manuscripts via multi-thousand-token prompt templates incurs high API costs, significant processing latency, and unpredictable reasoning hallucinations.

To resolve this dilemma, we design, implement, and empirically evaluate a \textbf{Relevance-Gated Agentic Workflow for Autonomous Research Monitoring}. Rather than adopting opaque, non-deterministic multi-agent frameworks (e.g., AutoGen~\cite{wu2023autogen}, CrewAI), our framework is anchored by an \textbf{explicit, deterministic finite-state machine}. Central to our architecture is a \textbf{three-way relevance gate} that partitions incoming publications into three distinct operational pathways:
\begin{enumerate}
    \item \textbf{HIGH Confidence Pathway ($s \ge t_{\text{high}}$):} Publications exhibiting high semantic alignment are directly forwarded to the user's research digest and processed by an interest-framed LLM summarizer.
    \item \textbf{BORDERLINE Pathway ($t_{\text{low}} < s < t_{\text{high}}$):} Ambiguous candidate publications are routed to an asynchronous human review queue, enabling researchers to triage borderline works with minimal manual effort.
    \item \textbf{LOW Confidence Pathway ($s \le t_{\text{low}}$):} Irrelevant publications are immediately pruned and logged into an auditable database table, preventing unnecessary downstream computation or LLM API calls.
\end{enumerate}

To prevent cold-start drift, relevance thresholds ($t_{\text{high}}, t_{\text{low}}$) are calibrated as percentiles ($P_{\text{high}}, P_{\text{low}}$) over a fixed, representative reference sample of 5,000 arXiv papers. As the user interacts with the system by triaging borderline papers and rating delivered summaries, the system accumulates multi-source feedback events. These signals trigger a dual-tier preference memory system: initially refining the interest vector using a Rocchio-style centroid adjustment~\cite{rocchio1971relevance}, and subsequently transitioning to a supervised, $L_2$-regularized logistic regression classifier once sufficient feedback data is accumulated and verified via out-of-fold generalization guards.

\subsection{Problem Statement and Motivation}
\label{sec:problem_statement}

Let $\mathcal{P} = \{p_1, p_2, \dots, p_N\}$ denote a continuous stream of incoming scientific papers, where each paper $p_i$ is characterized by its title, abstract, metadata, and dense embedding vector $\mathbf{p}_i \in \mathbb{R}^d$. A researcher's latent intellectual interests are captured by an evolving preference profile $\mathcal{U}$. The core operational challenge is to define a mapping function:
\begin{equation}
\Phi: \mathcal{P} \times \mathcal{U} \to \{\text{HIGH}, \text{BORDERLINE}, \text{LOW}\}
\end{equation}
such that:
\begin{enumerate}
    \item \textbf{Precision of High Route:} Publications auto-delivered to the digest maximize precision with respect to the user's primary interests, minimizing reading fatigue.
    \item \textbf{Safety of Low Route (Zero Discard of Relevant Papers):} The miss rate---the proportion of genuinely relevant papers erroneously assigned to LOW---is constrained to near zero ($\le 5\%$).
    \item \textbf{Economic and Computational Efficiency:} Generative LLM summarization is invoked exclusively on gate-cleared papers ($\text{HIGH} \cup \text{BORDERLINE}_{\text{accepted}}$), reducing token expenditures by over 90\% compared to unmitigated processing.
    \item \textbf{Active Personalization:} Human feedback on delivered digests and borderline review tasks progressively refines the decision boundary without requiring full model fine-tuning.
\end{enumerate}

\subsection{Research Gap}
\label{sec:research_gap}

Prior research in academic literature triage has remained bifurcated~\cite{tiropanis2020web}. On one hand, classical recommender systems (e.g., arXiv-Sanity~\cite{karpathy2015arxivsanity}) employ collaborative filtering or static TF-IDF/bag-of-words classifiers. While computationally inexpensive, they lack semantic nuance, struggle with lexical mismatch, and provide static ranking without adaptive agentic escalation. On the other hand, recent generative LLM agent workflows rely on multi-agent debate or comprehensive document reading, neglecting the computational and financial realities of high-throughput literature streams. Furthermore, existing human-in-the-loop systems frequently blur the distinction between reviewer triage decisions (evaluating unsummarized candidate abstracts) and consumer reactions (liking/disliking delivered syntheses), failing to isolate the informational value of different human interaction signals.

\subsection{Research Objectives and Contributions}
\label{sec:contributions}

In this work, we bridge this gap through five specific contributions:
\begin{enumerate}
    \item \textbf{Deterministic State-Machine Architecture:} We formulate an explicit Python state machine for autonomous research monitoring that enforces auditable transitions, idempotent cycle runs, and asynchronous human escalation without framework dependencies.
    \item \textbf{Calibrated Three-Way Percentile Gating:} We demonstrate that formulating relevance thresholds as percentiles of a fixed reference distribution provides robust, cold-start threshold stability that automatically transfers across shifting user profiles.
    \item \textbf{Informational Signal Ablation:} We evaluate the distinct learning values of end-user digest reactions versus human reviewer borderline triage decisions, proving that borderline review decisions accelerate preference adaptation significantly faster than high-confidence feedback alone.
    \item \textbf{Dual-Tier Preference Memory:} We implement and benchmark an adaptive progression that transitions from unsupervised Rocchio centroid drift to a supervised, regularized classifier governed by cross-validation promotion guards.
    \item \textbf{Rigorous Empirical and Boundary Reporting:} We ground all claims in an audited evaluation framework comprising 171 human-labeled papers, stratified Horvitz-Thompson sampling weights, frozen test access auditing, and an unvarnished analysis of live pilot constraints.
\end{enumerate}
